% ---
% id: meat-core/def:feasible
% kind: def
% title: "Feasible set and computable objective"
% summary: >-
%   The program minimises an entropy objective over the compiler's feasible set
%   ([[protocol-compiler/def:feasible]]), read as normalised block-diagonal
%   states meeting the declared marginal and constraints. Its infimum over the
%   acceptance set, h_low, is the value a key length is computed from.
% prominence: standard
% status: implemented
% uses: [def:hqkd-list, def:homog, foundations/def:kl, def:declaration, protocol-compiler/def:feasible]
% assumes: []
% refs: []
% provenance: newly-derived
% implements: certiqs_key.meat_core.objective.objective:check_feasible
% lean:
% notes: "The set, its conditions in both formats and its convexity and compactness are protocol-compiler/def:feasible's and are cited. The displays eq:feas, eq:feas-modes and eq:feas-general stand here under this contract's labels, in this module's reading of a block list as the block-diagonal operator it denotes, because the implementation and the tests name those labels and the implementation's refusal messages quote the line tags of eq:feas-general, the party-local line by the entry's data value party_local. The normalization paragraph, the objective and h_low are this module's own."
% ---

The set over which this module's program minimises is the compiler's
\tref{feasible set} (\xref{protocol-compiler/def:feasible}); the objective and
$\hlow$ are this module's alone.

\begin{definition}\label{def:feasible}
For a declaration $\Decl$, $\Feas$ is the \tref{feasible set} of
\xref{protocol-compiler/def:feasible}. Read as a set
of operators on $A\Qsys$, a list of sector blocks being the block-diagonal
operator it denotes, it is
\begin{equation}\label{eq:feas}
  \Feas\coloneqq\Big\{\!\omega\in\Pos{A\Qsys}:\
  (\idon{A}\otimes\Pi_s)\,\omega\,(\idon{A}\otimes\Pi_{s'})=0
  \ \forall s\ne s',\quad
  \tr\omega=1,\quad
  \omega\models\Decl,\quad
  \omega\models\CList\Big\},
\end{equation}
writing $\omega_s$ for the sector blocks, with $\omega\models\Decl$ the
marginal condition of \eqref{eq:margcond} and $\omega\models\CList$ the
constraint conditions in the form
Definition~\ref{def:homog} gives them for the declared mode. Written
out, $\Feas$ is
\begin{equation}\label{eq:feas-modes}
\begin{aligned}
  \text{trusted:}\quad&
  \omega_A=\tau_A,\quad
  \tr_{\Qsys}[(\idon{A}\otimes\hat O_k)\,\omega]\bowtie_k\Theta_k\
  \ \forall k;\\
  \text{untrusted:}\quad&
  \text{(no marginal condition)},\quad
  \tr[(\idon{A}\otimes\hat O_k)\,\omega]\bowtie_k c_k\ \ \forall k .
\end{aligned}
\end{equation}
The normalization $\tr\omega=1$ is written \emph{explicitly} in
\eqref{eq:feas}, and is not derived from anything: in
\tref{trusted-source mode} it also follows from $\omega_A=\tau_A$ (whence
$\tr\omega=\tr\tau_A=1$), but that derivation is unavailable in
\tref{untrusted-source mode}, where normalization is the \emph{only}
unconditional constraint the feasible set carries. Leaving it implied
would empty the untrusted-mode program of its one defining condition.
$\Feas$ is convex and compact in either mode and in either format
(\xref{protocol-compiler/def:feasible}). For
$\alpha\in(1,\infty)$,
$\mathbf q$ a distribution on $\Chatset$ and $\omega\in\Feas$, define the
\term[The function of an acceptance distribution and a feasible state that the program minimises: a divergence penalty plus one entropy term per key-live acceptance value.]{objective}
\begin{equation}\label{eq:objective}
  O_{\alpha}(\mathbf q,\omega)\coloneqq
  \frac{\alpha}{\alpha-1}\sum_{\hat c\notin\Chkey}
  q(\hat c)\log_2\frac{q(\hat c)}{P_{\hat c}(\omega)}
  +\sum_{\hat c\in\Chkey}\left[
  q(\hat c)\log_2\frac{q(\hat c)}{P_{\hat c}(\omega)}
  +\frac{q(\hat c)}{\alpha-1}\log_2\frac{q(\hat c)}{K_{\hat c}(\omega)}
  \right],
\end{equation}
with the conventions of \xref{foundations/def:kl} at vanishing arguments
($0\log_2\frac0v=0$; $u\log_2\frac u0=+\infty$ for $u>0$), and
\begin{equation}\label{eq:hlow-def}
  \hlow\coloneqq\inf_{\mathbf q\in\acSet}\ \inf_{\omega\in\Feas}\
  O_{\alpha}(\mathbf q,\omega).
\end{equation}

\emph{The general format.} For a declaration in the general format of
Definition~\ref{def:declaration}, $\Feas$ is the set of lists over the
\tref[joint sector]{joint sectors} of \xref{protocol-compiler/def:feasible},
$\omega=\bigoplus_{(t,s)}\omega^{(t,s)}$, with the $s$-block
$\hat O^{\uparrow(s)}_k$ of the \tref{lifted operator} of an entry scoped to $f_k$
and, for a party-local entry scoped to $f$, the split $s=(s_{f^c},s_f)$ and
the product $\Qsys_{f^c,s_{f^c}}$ of the other factors' blocks written
there, with every block positive semidefinite; its conditions read, block by
block,
\begin{equation}\label{eq:feas-general}
\begin{aligned}
  &\textstyle\sum_{(t,s)}\tr\omega^{(t,s)}=1
  &&\text{(normalization)};\\
  &\textstyle\sum_{s}\tr_{\Qsys_s}\omega^{(t,s)}=\tau^{(t)}_A
  \quad\forall t
  &&\text{(marginal, trusted-source mode)};\\
  &\textstyle\sum_{s}\tr_{\Qsys_s}\!\big[(\idon{A_t}\otimes
  \hat O^{\uparrow(s)}_k)\,\omega^{(t,s)}\big]\bowtie_k\Theta^{(t)}_k
  \quad\forall t
  &&\text{(trusted; }\Theta^{(t)}_k=c_k\tau^{(t)}_A
  \text{ if homogeneous)};\\
  &\textstyle\sum_{(t,s)}\tr\!\big[(\idon{A_t}\otimes
  \hat O^{\uparrow(s)}_k)\,\omega^{(t,s)}\big]\bowtie_kc_k
  &&\text{(scalar entry, untrusted-source mode)};\\
  &\textstyle\sum_{s_f\in\SecSet^{f}}\tr_{\Qsys_{f,s_f}}\!\Big[\big(\idon{}
  \otimes(\hat O^{(s_f)}_k-c_k\idon{})\big)\,
  \omega^{(t,(s_{f^c},s_f))}\Big]\bowtie_k0
  \quad\forall(t,s_{f^c})
  &&\text{(party-local, either mode)},
\end{aligned}
\end{equation}
the last as an operator on $\Hs{A_t}\otimes\Hs{\Qsys_{f^c,s_{f^c}}}$. The
objective \eqref{eq:objective} and $\hlow$ of \eqref{eq:hlow-def} are
defined on this $\Feas$ by the same formulas.
\end{definition}
